\documentclass[english, parskip=full]{scrartcl}
\usepackage[utf8]{inputenc}  % Kodierung der Datei
\usepackage[english]{babel}  % Sprache auf Deutsch 
\usepackage{color}
\usepackage{graphicx, gensymb}
\usepackage{amsfonts, cancel, amssymb}
\usepackage{amsmath, amsthm, array, esint}
\usepackage{booktabs, multirow}  % schönere Tabellen
\usepackage{caption}  % Anpassung der Captions
\usepackage{scrlayer-scrpage}    % Manipulation des Seitenstils
\pagestyle{scrheadings}  % Stil für die Seite setzen
\automark{section}  % Abschnittsnamen als Seitenbeschriftung 
\setheadsepline{.5pt}    % Kopzeile durch Linie abtrennen
\usepackage[hidelinks]{hyperref}  % Links und weitere PDF-Features
\usepackage{typearea, tikz, pgffor, verbatim}
\usetikzlibrary{calc, arrows.meta,arrows, graphs, quotes}
\usepackage[cache=false, outputdir=build]{minted}
\usemintedstyle{vs}
\usepackage{placeins} % provides \FloatBarrier

\definecolor{bg}{rgb}{0.9,0.9,0.9}
\newcommand\numberthis{\addtocounter{equation}{1}\tag{\theequation}}
\usepackage{svg}
\usepackage{siunitx}
\usepackage{physics}
\usepackage{tensor}
\renewcommand{\bra}{}
\newcommand{\kom}{}
\newcommand{\brak}{}
\renewcommand{\braket}{}
\newcommand{\intx}{}
\newcommand{\intr}{}
\newcommand{\kla}{}
\newcommand{\klam}{}
\newcommand{\klammer}{}
\newcommand{\oper}{}
\newcommand{\tecxt}{}
\newcommand{\Bra}{}
\newcommand{\Ket}{}
\renewcommand{\i}{\text{i}}
\newcommand{\e}{\text{e}}
\newcommand*{\Fourier}{\textsc{Fourier}\ }
\newcommand*{\F}{\mathcal{F}}
\newcommand*{\sinc}{\text{sinc\,}}
\newcommand{\conv}{\mathop{\scalebox{3.5}{\raisebox{-0.38ex}{\makebox[0.5\width][c]{$\ast$}}}}\limits}

\newcommand*{\titel}{2 Lattice polymer}
\newcommand*{\autor}{Joachim Schwardt}

\hypersetup{pdfauthor={\autor}, pdftitle={\titel}}

%\titlehead{titlehead}
\subject{Soft Condensed Matter}
\title{\titel}
\author{\autor}
\date{\today}

\begin{document}

%\begin{titlepage}
\maketitle  % Titel setzen
%\tableofcontents  % Inhaltsverzeichnis setzen
%\end{titlepage}

We consider a random walk on a one-dimensional lattice consisting of $N$ steps of size $\pm a$, where $a$ is the 1D lattice constant.
Starting at the origin, the walker reaches the position $R$ on the lattice.
The walk can be considered as a model for a one-dimensional ideal polymer with $N+1$ monomers having the end-to-end distance $R$ (in terms of lattice coordinates).
\\
The total number of configurations a polymer of length $N$ can have is simply given by $\Omega_0(N) = 2^N$, as there are two possible configurations for each (independent) step.
Let $N_\pm$ denote the number of steps in positive or negative direction respectively.
An end-to-end distance of $R$ is then equivalent to 
\begin{align}
R &= N_+ - N_- = N - 2N_-.
\end{align}
In this form it is evident that we are looking for the number of possible combinations of the $N_-$-steps, leading to
\begin{align}
\Omega(R,N) &= \binom{N}{N_-} = \binom{N}{\frac{N-R}{2}}.
\end{align}
The formula is not entirely correct, as it does not respect that e.g. odd $N$ and $R$ should give $\Omega(R,N)=0$.
This effectively means that only half the $R$-values are allowed with some non-zero probability. 
We will come back to this in the discussion of the next result.
\\
The probability for an end-to-end length is given by
\begin{align}
p(R) &= \frac{\Omega(R,N)}{\Omega_0(N)} = \frac{N!}{2^N}\frac{1}{\kla\left( \frac{N-R}{2} \right)! \kla\left( \frac{N+R}{2} \right)!}.
\end{align}
Consider the case $N\gg 1$ and $R\ll N$.
Then we may make use of the Stirling approximation $N!\sim\sqrt{2\pi N}(N/\e)^N$.
Note that $\log N! \sim N\log N - N + \frac{1}{2}\log (2\pi N)$.
It is easier to first consider the logarithm of the probability
\begin{align*}
\log p(R) &= -N\log 2 + \log N! - \log\kla\left( \frac{N-R}{2} \right)! - \log\kla\left( \frac{N+R}{2} \right)! \sim \\
&\sim -N\log2 + \log\sqrt{2\pi N} + N\log N - N - \log\kla\left( \pi \sqrt{N^2-R^2} \right) - \\
&\hspace{0.5cm} - \frac{N-R}{2}\log\frac{N-R}{2} + \frac{N-R}{2} - \frac{N+R}{2}\log\frac{N+R}{2} + \frac{N+R}{2} \simeq \\
&\simeq \log\kla\left( \frac{\sqrt{2\pi N}}{\pi N} \right) +\frac{R^2}{2N^2} + N\log \frac{N}{2} - \frac{N}{2}\log\kla\left( \frac{N^2}{4} \right) + \frac{R}{2}\log\frac{N-R}{N+R} = \numberthis \\
&= \frac{1}{2}\log \frac{2}{\pi N} +\frac{R^2}{2N^2} + \frac{R}{2} \klam\left[ \log\kla\left( 1 - \frac{R}{N} \right) - \log\kla\left( 1 + \frac{R}{N} \right) \right] \simeq \\
&\simeq \frac{1}{2}\log \frac{2}{\pi N} +\frac{R^2}{2N^2} + \frac{R}{2}\klam\left[ -\frac{R}{N} - \frac{R^2}{2N^2} - \frac{R}{N} + \frac{R^2}{2N^2} \right] = \\
&= \frac{1}{2}\log \frac{2}{\pi N} -\frac{R^2}{2N}.
\end{align*}
Thus we obtain the preliminary result
\begin{align}
p(R) &\sim \sqrt{\frac{2}{\pi N}} \exp\kla\left( -\frac{R^2}{2N} \right).
\end{align}
Recall that only half the $R$-values were allowed.
However, in this ``continuous'' limit this distinction is not very useful.
Instead, we just allow the entire range for $R$ but compensate by reducing the probability for any $R$ by two, giving
\begin{align}
p(R) &\sim \frac{1}{\sqrt{2\pi N}} \exp\kla\left( -\frac{R^2}{2N} \right).
\end{align}
In terms of the actual distance $r=Ra$ on the lattice we find
\begin{align}
p(r) &\sim \frac{1}{\sqrt{2\pi N}} \exp\kla\left( -\frac{r^2}{2Na^2} \right) \frac{\tecxt\text{d}R}{\tecxt\text{d}r} = \frac{1}{\sqrt{2\pi N}a} \exp\kla\left( -\frac{r^2}{2Na^2} \right).
\end{align}
The mean position is given by
\begin{align}
\bra\langle r \rangle &= \intr\int_{\mathbb{R}^{ }} rp(r) \, \mathrm{d}^{ }r = 0\quad (\tecxt\text{due to the symmetry }r\rightarrow -r).
\end{align}
The variance is then simply given by 
\begin{align}
\Delta r^2 &= \bra\langle r^2 \rangle = \intr\int_{\mathbb{R}^{ }} r^2p(r) \, \mathrm{d}^{ }r = \frac{1}{\sqrt{2\pi Na^2}} 2N^2a^2 \frac{\partial}{\partial N} \intr\int_{\mathbb{R}^{ }} \exp\kla\left( -\frac{r^2}{2Na^2} \right) \, \mathrm{d}^{ }r = \\
&= \frac{1}{\sqrt{2\pi Na^2}} 2N^2a^2 \frac{\partial}{\partial N} \sqrt{2\pi Na^2} = 2N^2a^2 \frac{1}{2N} = Na^2.
\end{align}
The entropy of the polymer for a given $r$ can be computed from $S(r) = k_B\log \Omega(r)$, where $\Omega(r) = p(r)\Omega$.
This leads to 
\begin{align}
S(r) &= k_B\log p(r)\Omega_0 = k_B\log\Omega + k_B\log p(r) =: S_0 - \frac{k_Br^2}{2\Delta r^2},
\end{align}
where $S_0 := k_B\log\Omega - \frac{1}{2}\log\kla\left( 2\pi\Delta r^2 \right)$ is independent of $r$.
Since $S''(r)<0$, entropy maximisation gives $r=0$ without any outer influence.
However, for a linear potential $U(r) = -\kappa r$ we find the free energy
\begin{align}
F(r) &= U(r) - TS(r) = -\kappa r - TS_0 + \frac{k_BTr^2}{2\Delta r^2}.
\end{align}
This function has a minimum at $F'(r) = 0 = -\kappa + \frac{k_BT}{\Delta r^2} \bar{r}$, giving $\bar{r}=\frac{\kappa \Delta r^2}{k_BT}$.
The assumed potential is equivalent to a constant force $-U'(r) = \kappa$.
For increasing temperature, the equilibrium length is $\bar{r} \propto \frac{\kappa}{T}$.
For this to be constant, the force would have to be directly proportional to temperature.
\\
The Central Limit Theorem would have directly provided us with the gaussian probability distribution, since the walker is the culmination of many independent identical events.

%This allows us to simplify the expression
%\begin{align}
%p(R) &\sim \frac{\sqrt{2\pi N} N^N\e^{-N}}{2^N} \frac{1}{\pi\sqrt{N^2-R^2}} \e^{\frac{N-R}{2}} \e^{\frac{N+R}{2}} \kla\left( \frac{N-R}{2} \right)^{-\frac{N-R}{2}} \kla\left( \frac{N+R}{2} \right)^{-\frac{N+R}{2}} \simeq \\
%&\simeq \frac{N^N}{2^{N-1}\sqrt{2\pi N}} \kla\left( \frac{N-R}{2} \right)^{-\frac{N-R}{2}} \kla\left( \frac{N+R}{2} \right)^{-\frac{N+R}{2}}
%\end{align}

\end{document}